% ============================================================
% Week 8 — Modeling
% Replace the empty \section{Results} placeholder.
%
% Upload to Overleaf:
%   fig10_baseline_confusion_matrix.pdf
% ============================================================

\section{Results}
\label{sec:results}

\subsection{Baseline Model}

As a reference, we trained a Logistic Regression classifier
(\texttt{class\_weight='balanced'}, $C=1.0$) on the TF-IDF features from
Section~III-G. Table~\ref{tab:baseline_results} shows the test-set
results (109~notifications).

\begin{table}[!t]
\centering
\caption{Baseline (TF-IDF + Logistic Regression) test-set results.}
\label{tab:baseline_results}
\begin{tabular}{lccc}
\hline
\textbf{Class} & \textbf{Prec.} & \textbf{Rec.} & \textbf{F1} \\
\hline
High   & 0.88 & 1.00 & 0.94 \\
Medium & 0.72 & 0.79 & 0.75 \\
Low    & 0.85 & 0.69 & 0.76 \\
\hline
\textbf{Macro avg} & 0.82 & 0.83 & 0.82 \\
\hline
\multicolumn{3}{l}{\textbf{Accuracy}} & 0.83 \\
\multicolumn{3}{l}{\textbf{Macro-F1}} & 0.8185 \\
\hline
\end{tabular}
\end{table}

The baseline achieves a macro-F1 of 0.8185 and 82.6\% accuracy.
High-urgency notifications are easiest to classify (F1\,=\,0.94) thanks
to distinctive keywords (\textit{minutes}, \textit{code}), while Medium
and Low are harder to separate (F1\,=\,0.75 and 0.76). The confusion
matrix (Fig.~\ref{fig:baseline_cm}) shows the main errors are Low
predicted as Medium (9~cases) and Medium predicted as Low (5~cases).

\begin{figure}[!t]
\centering
\includegraphics[width=0.85\columnwidth]{fig10_baseline_confusion_matrix.pdf}
\caption{Baseline confusion matrix on the test set.}
\label{fig:baseline_cm}
\end{figure}

\subsection{Transformer Models}

We fine-tune \texttt{bert-base-uncased}~\cite{devlin2019} and
\texttt{roberta-base}~\cite{liu2019} using
\texttt{AutoModelForSequenceClassification}~\cite{wolf2020}.
Table~\ref{tab:hyperparams} lists the hyperparameters.

\begin{table}[!t]
\centering
\caption{Transformer fine-tuning hyperparameters.}
\label{tab:hyperparams}
\begin{tabular}{ll}
\hline
\textbf{Parameter} & \textbf{Value} \\
\hline
Max sequence length     & 64 tokens \\
Batch size              & 16 \\
Learning rate           & $2 \times 10^{-5}$ \\
Optimiser               & AdamW (weight decay 0.01) \\
LR schedule             & Linear with 10\% warmup \\
Max epochs              & 5 \\
Early stopping patience & 2 (on val macro-F1) \\
Loss function           & Cross-entropy, class-weighted \\
Gradient clipping       & max norm = 1.0 \\
\hline
\end{tabular}
\end{table}

Class weights are set to $w_c = N / (3 \cdot n_c)$ to counteract the
mild imbalance. The best checkpoint (by val macro-F1) is saved and
re-evaluated on the test set, recording per-sample softmax probabilities
for the ranking evaluation. Training runs on the IE University HPC
cluster (NVIDIA RTX~6000 Ada, 48\,GB). BERT and RoBERTa results will be
added once the cluster job completes.
